\chapter{Head Lice (Pediculosis Capitis)}
\index{Head Lice (Pediculosis Capitis)}
\label{sec:Head Lice}

\section{Overview}


\begin{itemize}[noitemsep]
  \item Head lice is infestation of the scalp by \textit{Pediculus humanus capitis}, common in school-age children (6--12 million cases/year in the US). Transmission occurs through direct head-to-head contact
  \item Fomite transmission is less common.
\end{itemize}
\section{Causes}
This condition happens because of louse infestation with eggs (nits) laid at the hair base that hatch in 7--10 days and mature in 2--3 weeks.     
\section{Risk Factors}

\begin{itemize}[noitemsep]
  \item Genetic predisposition and family history
  \item Advancing age
  \item Lifestyle factors (diet, exercise, smoking, alcohol)
  \item Environmental exposures
  \item Underlying medical conditions
  \item Medication use
\end{itemize}
\section{Signs and Symptoms}

\subsection{Common symptoms}
\begin{itemize}[noitemsep]
  \item \item You may experience scalp itching (from hypersensitivity to louse saliva), visible nits (0.8 mm, yellow-white, firmly attached to the hair shaft), and less commonly live lice. Pain or discomfort in the affected area    \item Pain or discomfort in the affected area
  \end{itemize}

\subsection{Emergency warning signs}
\begin{itemize}[noitemsep]
  \item Severe or worsening symptoms of head lice (pediculosis capitis) require immediate medical evaluation
\end{itemize}
\section{Progression and Stages}
The condition may progress through different stages of severity over time. Early stages often involve mild or subtle symptoms that may be overlooked. Moderate stages involve more noticeable symptoms affecting daily function. Advanced stages may involve significant impairment and complications.
\section{Complications}
\begin{itemize}[noitemsep]
  \item Disease progression leading to worsening symptoms
  \item Secondary infections or injuries
  \item Side effects of treatment
  \item Reduced quality of life
  \item Emotional and psychological impact
\end{itemize}
\section{Diagnosis}
Doctors diagnose this using visual inspection with a fine-toothed comb

\begin{itemize}[noitemsep]
  \item Wood's lamp shows nits as pale blue. The first-choice treatment typically involves permethrin 1\% or pyrethrins (OTC), applied for 10 minutes and repeated in 7--10 days
  \item Second-line options include malathion 0.5\% lotion, benzyl alcohol 5\%, spinosad 0.9\%, and ivermectin 0.5\% lotion
  \item Oral ivermectin is used for refractory cases. Wet combing is a non-chemical option. All close contacts should be checked
  \item Environmental measures include washing bedding in hot water.
  \item Comprehensive medical history and physical examination
  \item Laboratory tests to assess organ function
  \item Imaging studies to visualize affected structures
  \item Specialized testing depending on the affected system
  \item Consultation with relevant specialists
\end{itemize}
\section{Prevention}
\begin{itemize}[noitemsep]
  \item Maintain a healthy lifestyle with balanced diet and exercise
  \item Avoid known risk factors
  \item Attend regular medical check-ups for early detection
  \item Follow recommended screening guidelines
\end{itemize}
\section{Treatment Options}
Medications to manage symptoms and address underlying mechanisms Lifestyle modifications and self-care measures Physical therapy and rehabilitation Surgical intervention when indicated Regular monitoring and follow-up care Patient education and support services

\begin{itemize}[noitemsep]
  \item Medications to manage symptoms and address underlying mechanisms
  \item Lifestyle modifications and self-care measures
  \item Physical therapy and rehabilitation
  \item Surgical intervention when indicated
  \item Regular monitoring and follow-up care
\end{itemize}
\section{Living with It}
\begin{itemize}[noitemsep]
  \item Follow your treatment plan as prescribed
  \item Maintain regular follow-up with your healthcare provider
  \item Make necessary lifestyle adjustments
  \item Seek support from family, friends, or support groups
  \item Stay informed about your condition
\end{itemize}
\section{Outlook}
\begin{itemize}[noitemsep]
  \item Most people recover well with appropriate treatment
  \item No-nit policies in schools are not recommended by the AAP/CDC.
\end{itemize}
